% Part: first-order-logic
% Chapter: introduction
% Section: sentences

\documentclass[../../../include/open-logic-section]{subfiles}

\begin{document}

\olfileid{fol}{int}{snt}

\olsection{语句}

现在，我们（初步）定义了结构与公式之间的满足关系（即“在……中为真”）。但这一定义需要一个额外的要素——变元指派——而我们想要的是语句的定义。我们该如何去掉变元指派？语句又是什么？

你可能还记得初学形式逻辑时，关于变元与量词之间关系的讨论。量词后总是跟着一个变元，而在该量词所作用的句子部分（即其“辖域”）内，该变元就被理解为被那个量词所“约束”。在公式中，并不要求每个变元都有与之匹配的量词，没有匹配量词的变元就是“自由的”或“不受约束的”。我们将把所有没有自由变元的公式称为语句。

同样，变元在公式~$!A$中的某次出现何时被约束、被哪个量词约束、何时是自由的，这些直观概念并不难把握。你可能学过一种检验方法，可能需要数括号。我们必须给出一个精确的定义——而且由于公式是归纳定义的，我们也可以归纳地定义变元~$x$在公式~$!A$中的自由出现与约束出现。例如，对于我们简化的语言，它可以这样给出：

\begin{enumerate}
  \item 如果 $!A$ 是原子公式，则 $x$ 在其中的所有出现都是自由的（也就是说，$x$ 在~$\Atom{\Obj P}{x}$ 中的出现是自由的）。
  \item 如果 $!A$ 具有形式 $\lnot !B$，则 $\lnot !B$ 中 $x$ 的一次出现是自由的，当且仅当 $x$ 在 $!B$ 中的对应出现是自由的（也就是说，$!B$ 中变元的自由出现恰好就是 $\lnot !B$ 中的对应出现）。
  \item 如果 $!A$ 具有形式 $(!B \land !C)$，则 $(!B \land !C)$ 中 $x$ 的一次出现是自由的，当且仅当 $x$ 在 $!B$ 或 $!C$ 中的对应出现是自由的。
  \item 如果 $!A$ 具有形式 $\lexists[x][!B]$，则 $!A$ 中 $x$ 的任何出现都不是自由的；如果 $!A$ 具有形式 $\lexists[y][!B]$，其中 $y$ 是与 $x$ 不同的变元，则 $\lexists[y][!B]$ 中 $x$ 的一次出现是自由的，当且仅当 $x$ 在 $!B$ 中的对应出现是自由的。
\end{enumerate}

一旦我们有了变元的自由出现与约束出现的精确定义，我们就可以简单地说：语句就是没有任何自由变元出现的公式。

\end{document}